\documentclass[11pt]{article}
\usepackage{cse216,hyperref}
\def\title{HW 4}
\begin{document}
\maketitle

\section*{Due October 27 at 3:00 pm}

\textbf{(2 questions, 110 points total)}

\begin{qunlist}

  \q{40}{EUF}\\
  Use the congruence closure algorithm to decide which of the following EUF conjunctions is satisfiable.
  Show the steps of the algorithm.
  \begin{qparts}
  \item
    $f(x) = x \land f(g(x)) \ne g(x)$

    \begin{answer}
      the initial partition has 4 terms
      \begin{equation}
        \{\{f(x)\},\{x\},\{f(g(x))\},\{g(x)\}\}
      \end{equation}

      there is only one new equivalence class

      \begin{equation}
        \{\{f(x),x\},\{f(g(x))\},\{g(x)\}\}
      \end{equation}

      none of the non equalities cause a contradiction with the equivalence class created by the equalities thus the EUF conjunction is sat.

    \end{answer}

  \item
    $f(x,y) = f(y,x) \land g(y) \ne g(f(y, g(y))) \land f(g(y), y) = y \land x = g(y)$

    \begin{answer}
      the subterms that show up are
      \begin{equation}
        \{\{x\}, \{y\}, \{g(y)\}, \{f(x,y)\}, \{f(y,x)\}, \{f(y, g(y))\}, \{f(g(y), y)\}, \{g(f(y, g(y)))\}\}
      \end{equation}
      the equalites asserted makes

      \begin{equation}
        \{\{x, g(y)\}, \{y, f(g(y), y)\}, \{f(x,y), f(y,x)\}, \{f(y, g(y))\}, \{g(f(y, g(y)))\}\}
      \end{equation}

      note  $x \sim g(y)$ and $y \sim y$.

      thus
      \begin{equation}
        \{\{x, g(y)\}, \{y, f(g(y), y), f(x,y), f(y,x)\}, \{f(y, g(y))\}, \{g(f(y, g(y)))\}\}
      \end{equation}


      also
      $y \sim y$ and $x \sim g(y)$.

      which gives

      \begin{equation}
        \{\{x, g(y)\}, \{y, f(g(y), y), f(x,y), f(y,x), f(y, g(y))\}, \{g(f(y, g(y)))\}\}
      \end{equation}

      finally since
      $y \sim f(y, g(y))$

      \begin{equation}
        \{\{x, g(y), g(f(y, g(y)))\}, \{y, f(g(y), y), f(x,y), f(y,x), f(y, g(y))\}\}
      \end{equation}
      to check inequalities

      $g(y)$ is in class $\{x, g(y), g(f(y, g(y)))\}$ but $g(f(y, g(y)))$ is in the same class, therefore there is a contradiction.
    \end{answer}
  \end{qparts}

  \q{70}{Lazy SMT solving}\\
  Consider the following Linear Integer Arithmetic (LIA) formula $\phi_{\text{SMT}}$:
  \begin{align*}
    &(x > 5 \;\lor\; \lnot (y > 3)) \\
    \land & (x \le 4 \;\lor\; x + y < 9) \\
    \land & (y > 3 \;\lor\; y > x) \\
    \land & (x > 3)
  \end{align*}
  We'll use the lazy SMT algorithm to solve this.
  You can play the part of the LIA conjunction decision procedure yourself (e.g.~you can tell that $x > 10 \land x < 3$ is UNSAT, since no integer is simultaneously greater than 10 and less than 3).
  You don't need to show the internals of the SAT solving: feel free to use any SAT solver, e.g.~the online version of CryptoMiniSat \href{https://msoos.github.io/cryptominisat_web/}{here}.
  \begin{qparts}
  \item
    What is the Boolean abstraction $\phi_{\text{SAT}}$ of $\phi_{\text{SMT}}$?

    \begin{answer}

      let
      \begin{itemize}
        \item $A := (x > 5)$

        \item $B := (y > 3)$

        \item $C := (x \le 4)$

        \item $D := (x + y < 9)$

        \item $E := (y > x)$

        \item $F := (x > 3)$
      \end{itemize}

      we get

      \begin{equation}
        \phi_{\text{SAT}} = (A \lor \lnot B) \land (C \lor D) \land (B \lor E) \land F
      \end{equation}

    \end{answer}

  \item
    Give a satisfying assignment of $\phi_{\text{SAT}}$, and the corresponding conjunction of LIA literals.

    \begin{answer}

      let
      \begin{equation}
        A = \top, \quad B = \top, \quad C = \top, \quad D = \bot, \quad E = \bot, \quad F = \top
      \end{equation}

      which is the same as

      \begin{equation}
        (x > 5) \land (y > 3) \land (x \le 4) \land (x + y \ge 9) \land (y \le x) \land (x > 3)
      \end{equation}

    \end{answer}

  \item
    Is this conjunction satisfiable?
    If so, give a satisfying assignment.
    If not, give a theory lemma which has as few literals as possible.

    \begin{answer}
      no it is not. x cant be greater than 5 and less than 4.

      new theory lemma:
      \begin{equation}
        \lnot (x > 5) \lor \lnot (x \le 4)
      \end{equation}

    \end{answer}

  \item
    Continue the lazy SMT algorithm until you find a satisfying assignment for $\phi_{\text{SMT}}$ or prove it unsatisfiable.
    Show each solution you find to the Boolean abstraction and each theory lemma you add.

    \begin{answer}
      \begin{equation}
        A = \top, \quad B = \top, \quad C = \bot, \quad D = \top, \quad E = \bot, \quad F = \top
      \end{equation}

      which gives
      \begin{equation}
        (x > 5) \land (y > 3) \land (x > 4) \land (x + y < 9) \land (y \le x) \land (x > 3)
      \end{equation}

      but this also breaks because of

      \begin{equation}
        \lnot (x > 5) \lor \lnot (y > 3) \lor \lnot (x + y < 9)
      \end{equation}

      \begin{equation}
        A = \bot, \quad B = \bot, \quad C = \bot, \quad D = \top, \quad E = \top, \quad F = \top
      \end{equation}

      \begin{equation}
        (x \le 5) \land (y \le 3) \land (x > 4) \land (x + y < 9) \land (y > x) \land (x > 3)
      \end{equation}

      which breaks
      \begin{equation}
        (x \le 4) \lor (y > 3) \lor \lnot (y > x) \quad (\text{or } C \lor B \lor \lnot E)
      \end{equation}

      \begin{equation}
        A = \bot, \quad B = \bot, \quad C = \top, \quad D = \top, \quad E = \top, \quad F = \top
      \end{equation}

      \begin{equation}
        (x \le 5) \land (y \le 3) \land (x \le 4) \land (x + y < 9) \land (y > x) \land (x > 3)
      \end{equation}
      which adds
      \begin{equation}
        \lnot (x \le 4) \lor \lnot (y \le 3) \lor \lnot (y > x)
      \end{equation}

      but the final formula with all lemmas is unsat so the original expression is unsat.
    \end{answer}
  \end{qparts}

\end{qunlist}
\end{document}
